\documentclass[11pt]{article}
\usepackage[a4paper,margin=1in,bottom=1.18in]{geometry}
\usepackage{amsmath,amssymb,amsthm,mathtools}
\usepackage{graphicx}
\usepackage{microtype}
\usepackage[hidelinks]{hyperref}

\newtheorem{theorem}{Theorem}
\newtheorem{lemma}[theorem]{Lemma}
\newtheorem{corollary}[theorem]{Corollary}
\newcommand{\Irr}{\operatorname{Irr}}
\newcommand{\cp}{\operatorname{cp}}
\newcommand{\cd}{\operatorname{cd}}

\title{Finite Support of Projected Plancherel Weak Limits\\
\large A partial answer to Question (I) in arXiv:1712.04772}
\author{Automated mathematics run \texttt{fa\_banach\_001}}
\date{9 August 2026}

\begin{document}
\maketitle

\noindent\fbox{\parbox{0.94\textwidth}{\textbf{Status.}
Candidate substantial partial result, likely valid.  The theorem rules out
all infinite-support limits.  It does not classify the finitely supported
limits.}}

\section{Source question}

The source is R.~Adamczak, \emph{Random non-Abelian $G$-circulant
matrices. Spectrum of random convolution operators on large finite groups},
arXiv:1712.04772 \cite{Adamczak}.  Section~7, Question~(I), page~25 asks for a
characterization of the weak limits of projected Plancherel measures and notes
that the examples in the paper contain only finite mixtures.

\begin{figure}[ht]
  \centering
  \includegraphics[width=\textwidth]{figures/open_problem_crop.png}
  \caption{The complete source passage containing Question~(I), rendered from
  page~25 of the source paper.}
\end{figure}

For a finite group $G$, write $\Irr(G)$ for its irreducible complex
characters and put
\[
 \mu_G=\sum_{\chi\in\Irr(G)}\frac{\chi(1)^2}{|G|}\,\delta_{\chi(1)}.
\]
This is a probability measure on
$\overline{\mathbb N}=\mathbb N\cup\{\infty\}$, where
$\overline{\mathbb N}$ has its one-point compactification topology.

\section{Result and proof intuition}

\subsection*{Proof intuition}

The projected Plancherel measure remembers the commuting probability through
one continuous moment:
\[
 \cp(G)=\frac{|\Irr(G)|}{|G|}
       =\int_{\overline{\mathbb N}}d^{-2}\,d\mu_G(d),
 \qquad \infty^{-2}:=0.
\]
Consequently, a limiting measure with any mass at finite integers forces a
uniform positive lower bound on commuting probability.  P.~M.~Neumann's
theorem then puts the groups, uniformly, within bounded index of a group that
is finite-by-abelian with bounded finite kernel.

The remaining point is not that the character degrees themselves are bounded
--- extraspecial groups show this is false --- but that their \emph{number of
distinct values} is bounded.  Clifford theory reduces each degree to a
constituent over the bounded finite kernel and a projective representation of
an abelian quotient.  For a fixed multiplier, all irreducible projective
representations of a finite abelian group have the same degree.  Thus only
boundedly many degree values occur.  Support cardinality at most $D$ is closed
under weak convergence, which finishes the argument.

\begin{theorem}\label{thm:main}
Let $(G_j)$ be finite groups with $|G_j|\to\infty$.  If
$\mu_{G_j}$ converges weakly on $\overline{\mathbb N}$ to a probability
measure $\mu$, then $\mu$ has finite support.

More quantitatively, if
\[
 c:=\int_{\overline{\mathbb N}}d^{-2}\,d\mu(d)>0,
\]
then $|\operatorname{supp}\mu|$ is bounded by a constant depending only on
$c$.
\end{theorem}

Thus the infinite-support examples suggested in Question~(I) do not exist.
The exact characterization of the possible finite-support measures remains
open.

\section{The finite-by-abelian degree lemma}

For a finite group $K$, let
$\cd(K)=\{\chi(1):\chi\in\Irr(K)\}$.

\begin{lemma}\label{lem:degrees}
For every pair of positive integers $a,b$ there is $D(a,b)$ with the
following property.  Suppose $G$ is finite and has normal subgroups
$H\leq \Gamma\leq G$ such that
\[
 [G:\Gamma]\leq a,\qquad |H|\leq b,
 \qquad\text{and}\qquad \Gamma/H\ \text{is abelian}.
\]
Then $|\cd(G)|\leq D(a,b)$.
\end{lemma}

\begin{proof}
First consider $\Gamma$.  Given $\eta\in\Irr(\Gamma)$, choose an irreducible
constituent $\theta\in\Irr(H)$ of $\eta|_H$, and let
$I=I_\Gamma(\theta)$ be its inertia group.  Clifford correspondence gives
\[
 \eta=\operatorname{Ind}_I^\Gamma\psi,
 \qquad
 \eta(1)=[\Gamma:I]\,\psi(1)
\]
for some $\psi\in\Irr(I)$ lying over $\theta$.

The character $\theta$ is $I$-invariant.  The character-triple form of
Clifford theory identifies the characters of $I$ lying over $\theta$ with
irreducible projective representations of the abelian group $I/H$ for one
fixed multiplier, with all degrees multiplied by $\theta(1)$.  For a finite
abelian group $A$ and a fixed multiplier $\alpha$, every irreducible
$\alpha$-projective representation has the same degree.  Indeed, if $R$ is
the radical of the commutator bicharacter of $\alpha$, the twisted group
algebra $\mathbb C^\alpha[A]$ is a direct sum of equal matrix blocks of size
$\sqrt{|A|/|R|}$.

\pagebreak
It follows that, for each $\Gamma$-orbit in $\Irr(H)$, all irreducible
characters of $\Gamma$ lying over that orbit have the same degree.  Hence
\[
 |\cd(\Gamma)|\leq |\Irr(H)|\leq |H|\leq b.
\]

Now take $\chi\in\Irr(G)$ and a constituent
$\eta\in\Irr(\Gamma)$ of $\chi|_\Gamma$.  A second application of Clifford
theory gives
\[
 \chi(1)=t e\,\eta(1),
\]
where $t$ is the size of the $G$-orbit of $\eta$ and $e$ is the degree of an
irreducible projective representation of a subgroup of $G/\Gamma$.  In
particular, $1\leq t\leq a$ and $1\leq e\leq a$.  Thus each of the at most
$b$ values of $\eta(1)$ gives at most $a^2$ values of $\chi(1)$.  One may
therefore take $D(a,b)=a^2b$.
\end{proof}

\section{Proof of the main theorem}

\begin{proof}[Proof of Theorem~\ref{thm:main}]
Define $f:\overline{\mathbb N}\to\mathbb R$ by
$f(d)=d^{-2}$ for $d\in\mathbb N$ and $f(\infty)=0$.  This is continuous.
The character-degree sum formula and the equality between the number of
irreducible characters and the number of conjugacy classes give
\[
 \int f\,d\mu_G
 =\sum_{\chi\in\Irr(G)}\frac{\chi(1)^2}{|G|}\frac1{\chi(1)^2}
 =\frac{|\Irr(G)|}{|G|}
 =\cp(G).
\]
Therefore weak convergence implies $\cp(G_j)\to c$.

If $c=0$, positivity of $f$ at every finite integer forces
$\mu=\delta_\infty$, so the conclusion holds.  Suppose $c>0$.  For all large
$j$, $\cp(G_j)\geq c/2$.  P.~M.~Neumann's theorem, in the quantitative form
quoted in \cite[Theorem~1.2]{Tointon}, supplies constants $a=a(c)$ and
$b=b(c)$ and normal subgroups
\[
 H_j\leq\Gamma_j\leq G_j
\]
such that $[G_j:\Gamma_j]\leq a$, $|H_j|\leq b$, and
$\Gamma_j/H_j$ is abelian.  Lemma~\ref{lem:degrees} shows that every
$\mu_{G_j}$ is supported on at most $D(a,b)$ points.

It remains only to pass this bound to the weak limit.  If $\mu$ had
$D(a,b)+1$ distinct support points, choose pairwise disjoint open
neighborhoods of those points.  Each has positive $\mu$-measure.  By the
Portmanteau theorem, all these neighborhoods have positive
$\mu_{G_j}$-measure for all sufficiently large $j$, contradicting the support
bound.  Hence $\mu$ has finite support.
\end{proof}

\section{A mixed finite--infinite example}

The conclusion cannot be strengthened to a dichotomy between finite support
inside $\mathbb N$ and $\delta_\infty$.  Fix a prime $p$ and let $E_m$ be an
extraspecial $p$-group of order $p^{1+2m}$.  It has $p^{2m}$ linear
characters and $p-1$ irreducible characters of degree $p^m$.  Therefore
\[
 \mu_{E_m}=\frac1p\delta_1+\frac{p-1}{p}\delta_{p^m}
 \ \Longrightarrow\ 
 \frac1p\delta_1+\frac{p-1}{p}\delta_\infty.
\]
This also explains why bounded commuting probability does not bound the
degrees themselves; the bounded quantity needed in the proof is the number
of distinct degree values.

\section{Scope, novelty, and verification}

The theorem settles the infinite-support part of Question~(I), but it does not
determine which finitely supported measures occur.  In particular, arithmetic
and representation-theoretic restrictions on their atoms and weights remain
to be classified.

A bounded search covered the run's solution, attempt, proof-gap, and registry
indexes; the local parsed arXiv corpus; and targeted primary-source searches
for projected Plancherel limits, character-degree distributions, and commuting
probability.  The search found Neumann's structure theorem and work on scalar
commuting-probability limits, but no explicit statement of
Theorem~\ref{thm:main}.  Novelty is plausible, not certified.

The proof is non-computational.  Human review should focus on
Lemma~\ref{lem:degrees}, especially the standard character-triple reduction
and the equal-degree decomposition of a twisted group algebra of a finite
abelian group.  The source crop and every final packet page were rendered and
visually inspected.

\begin{thebibliography}{9}

\bibitem{Adamczak}
R.~Adamczak,
\emph{Random non-Abelian $G$-circulant matrices. Spectrum of random
convolution operators on large finite groups},
Electron. J. Probab. \textbf{24} (2019), paper no.~42;
arXiv:1712.04772.

\bibitem{Tointon}
M.~Tointon,
\emph{Commuting probability of infinite groups},
J. Lond. Math. Soc. (2) \textbf{101} (2020), 1280--1297;
arXiv:1707.05565.  See Theorem~1.2 for the finite-group theorem of
P.~M.~Neumann used here.

\bibitem{Isaacs}
I.~M.~Isaacs,
\emph{Character Theory of Finite Groups},
Dover Publications, 1994.

\bibitem{Eberhard}
S.~Eberhard,
\emph{Commuting probabilities of finite groups},
Bull. Lond. Math. Soc. \textbf{47} (2015), 796--808;
arXiv:1411.0848.

\end{thebibliography}

\end{document}
